\section{Security}\label{sec:security}

\subsection{Threats considered}
The application targets a closed LAN, so the main threats are other hosts on the same network:
reading traffic with a packet sniffer, modifying or injecting packets, impersonating a peer at the
tracker, and guessing passwords. \Cref{tab:threats} maps each threat to its countermeasure.

\begin{table}[H]\centering\footnotesize
\caption{Threats and countermeasures}\label{tab:threats}
\begin{tabularx}{\textwidth}{@{}>{\raggedright\arraybackslash}p{4.0cm} Y@{}}
\toprule
\textbf{Threat} & \textbf{Countermeasure} \\
\midrule
Reading chat or files on the wire & AES-256-GCM on every chat text and every file chunk, random 96-bit nonce per message \cite{nist80038d} \\
Changing sender, channel or recipient of a message & these fields are AAD of AES-GCM: any change invalidates the tag \\
Moving a chunk to another position or transfer & AAD of a chunk is \code{"transfer_id:seq"} \\
Forging \code{START} or \code{ACK} (e.g.\ a fake ``complete'') & HMAC-SHA256 over the canonical JSON \cite{rfc2104} \\
Corrupted file & SHA-256 of the whole file checked before the file is kept \\
Taking over another peer's registration & per-peer token, required for every change; constant-time comparison \\
Stolen password database & PBKDF2-HMAC-SHA256, 200\,000 iterations, 16-byte random salt per user \cite{rfc8018} \\
Session theft from JavaScript & cookie \code{HttpOnly; SameSite=Lax}, value signed with HMAC, server-side session list \\
Path traversal (\code{../}) & received file names reduced to their base name; static files checked to stay inside the document root \\
Flooding with huge requests & limits: \SI{64}{KiB} headers and frames, \SI{100}{MiB} files, 2000-character messages \\
\bottomrule
\end{tabularx}
\end{table}

\subsection{Keys}
All peers of a group share one passphrase (\code{--shared-key}). From it, HMAC-SHA256 with distinct
labels derives three independent keys: one for chat texts, one for file chunks, one for UDP control
packets. Using a key for exactly one purpose means that, for example, a chat ciphertext can never be
accepted as a file chunk.

\pysnippet[lst:encrypt]{encrypt_chat.py}{Chat encryption with the visible fields as associated data}

\subsection{Authentication}
Passwords are checked only by the tracker, which also learns which user signs in on which peer --- the
``server must track connected users'' of rubric step~3. The peer daemon then issues its own session
cookie to the browser (RFC~6265 \cite{rfc6265}). Each peer uses its own cookie name, because browsers
share cookies between ports of the same host. A user whose password is wrong gets the same answer
whether or not the user name exists, and both cases take the same time (a dummy PBKDF2 run).

The tracker also lets receivers check the \code{sender} field: the shared key proves that a frame
comes from a member of the group, and the tracker's record says which user is signed in on the
origin peer. When the two disagree, the interface marks the message ``unverified''.

\subsection{Limitations}
The design is honest about what it does not protect. Everybody who holds the shared key can read all
messages and can, in principle, forge a frame with any \code{sender}; the ``unverified'' mark helps,
but per-user keys (for example a key exchange with the tracker as certificate authority) would be the
real fix. HTTP between browser, peer and tracker is not encrypted, so on an untrusted network the
login request would be visible; TLS on those connections would remove that. The demo accounts share
one password, which is acceptable for a demonstration only.
